\setcounter{definition}{0}
\setcounter{theorem}{0}
\subsection{Gyökei}
\subsubsection{Fogalma}
\begin{definition}
	Egy $f$ polinomnak az $x_0$ érték a gyöke,
	ha a megfelel polinomfüggvény ott a $0$ értéket veszi fel.
	\begin{equation*}
		x_0 \overset{f}{\mapsto} f(x_0) = 0
	\end{equation*}
\end{definition}
\subsubsection{Fokok és gyökszám}
\begin{theorem}
	Legyen $\mathbb{K} = \{ \mathbb{C, R, Q, Z}_p \}$.
	Egy $f \in \mathbb{K}[x]$ polinomnak legfeljebb $\deg f$ gyöke lehet.
\end{theorem}
\begin{remark*}
	A tétel $\mathbb{Z}_8[x]$ nem igaz.
\end{remark*}
\begin{proof}
	\begin{gather*}
		\deg f = 0 \implies f = c_0, c_0 \neq 0 \implies
		f\text{-nek nincs gyöke} \\
		\deg f \geq 1 \land f\text{-nek nincs gyöke} \implies
		f\text{-nek nincs gyöke} \\
		\deg f \geq 1 \land f(x_0) = 0 \implies
		\exists q, r \in \mathbb{K}[x]: f = (x - x_0) \cdot q + r \quad
		\deg r < 1
	\end{gather*}
	\begin{multline*}
		f(x_0) = 0 = q(x_0) \cdot (x_0 - x_0) + r \implies r = 0 \implies \\
		f = (x - x_0) \cdot q \implies \deg q = \deg f - 1
	\end{multline*}
	\begin{equation*}
		\exists x_1: x_1 \neq x_0 \land f(x_1) = 0 \implies
		f(x_1) = (x_1 - x_0) \cdot q(x_1) \implies q(x_1) = 0
	\end{equation*}
	Mivel $q$-nak legfeljebb $n - 1$ gyöke lehet,
	így $f$-nek legfeljebb $n$ gyöke lehet.
\end{proof}
